% Critère d'évolution spontanée : position du quotient de réaction par rapport à la constante d'équilibre (L'essentiel, chapitre 7)
\begin{tikzpicture}[x=1mm,y=1mm]
\draw[otMuted,line width=0.6pt,-{Stealth[length=2.2mm]}] (0,0) -- (130,0) node[right,pcLegende] {$Q_\mathrm{r}$};
\draw[otInk,line width=0.8pt] (65,-2.5) -- (65,2.5) node[above,font=\otMedium\fontsize{9}{10}\selectfont,text=otInk] {$K$};
\fill[otGreenTint] (2,-14) rectangle (62,-5);
\fill[otNavyTint] (68,-14) rectangle (128,-5);
\node[pcLegendeG,text=otGreenDark,align=center] at (32,-9.5) {$Q_\mathrm{r}<K$ : évolution dans le sens direct};
\node[pcLegendeG,text=otNavy,align=center] at (98,-9.5) {$Q_\mathrm{r}>K$ : évolution dans le sens inverse};
\draw[otGreenDark,line width=1.2pt,-{Stealth[length=2.6mm]}] (20,5) -- (58,5);
\draw[otNavy,line width=1.2pt,-{Stealth[length=2.6mm]}] (110,5) -- (72,5);
\node[pcLegende,text=otGreenDark,anchor=south] at (39,6) {$Q_\mathrm{r}$ augmente};
\node[pcLegende,text=otNavy,anchor=south] at (91,6) {$Q_\mathrm{r}$ diminue};
\node[pcLegende,anchor=north,align=center] at (65,-15.5) {$Q_\mathrm{r}=K$ : état d'équilibre,\\pas d'évolution macroscopique};
\end{tikzpicture}
\par\smallskip\legende{Un système hors équilibre évolue spontanément de façon que son quotient de réaction se rapproche de $K$.}
